\section{Conclusion}
In this work, we explore the role of higher-dimensional topological structures in retrieval-augmented graph learning. We propose a novel framework, namely \underline{Re}trieved Cellular \underline{T}opologies-\underline{A}ugmented \underline{G}raph Learning (ReTAG), which lifts input graphs into cellular complexes and constructs a knowledge base of multi-dimensional topology-aware subgraphs, termed cellular topologies. During inference, ReTAG retrieves structurally aligned cellular topologies and incorporates them via a multi-dimensional message-passing mechanism that captures complex topological dependencies beyond conventional pairwise relations. Furthermore, a cellular topological contrastive learning module is introduced to reinforce high-dimensional structural semantics by aligning features within individual topological cells. Extensive experiments across multiple graph tasks demonstrate that ReTAG outperforms state-of-the-art methods.